\reportheading{5. Comparison with lighting requirements}

Table 3 compares the DIALux results with EN 12464-1:2021 and the lighting power limit in BEC 2024. The office activity and room assumptions are described in Section 1.

\tabletitle{Table 3. Adopted requirements and calculated performance.}

\begin{reporttable}{L{0.3}L{0.34}L{0.36}}
\rowcolor{tablegrey} \textbf{Quantity} & \textbf{Requirement} & \textbf{Calculated result} \\ \hline
Desk mean & ≥500 lx & Lowest mean: 593 lx \\
Desk uniformity & \(U_{\mathrm{o}}\) ≥0.60 & Lowest \(U_{\mathrm{o}}\): 0.80 \\
Directional glare & UGR ≤19 & Highest UGR: 18.1 \\
Colour rendering & Ra ≥80 & Ra >80 (manufacturer) \\
Seated cylindrical & ≥150 lx; \(U_{\mathrm{o}}\) ≥0.10 & 224 lx; \(U_{\mathrm{o}}\) ≥0.59* \\
Immediate surround & ≥300 lx; \(U_{\mathrm{o}}\) ≥0.40 & 595 lx; \(U_{\mathrm{o}}\) ≥0.64* \\
Background floor & ≥100 lx; \(U_{\mathrm{o}}\) ≥0.10 & 617 lx; \(U_{\mathrm{o}}\) =0.69 \\
Exposed walls & ≥150 lx; \(U_{\mathrm{o}}\) ≥0.10 & 216 lx; \(U_{\mathrm{o}}\) ≥0.24* \\
Exposed ceiling & ≥100 lx; \(U_{\mathrm{o}}\) ≥0.10 & 137 lx; \(U_{\mathrm{o}}\) =0.89 \\
Lighting power density & ≤7.2 W/m² & 4.153 W/m² including allowance \\
Manual control points & ≥32 for 1042.10 m² & 32 proposed \\
Occupancy shutoff & ≤15 min & 10 min specified \\
Daylight response & Conditional on glazing and load & Pending glazing survey \\
\end{reporttable}

\tablenote{Sources: [1, Table 34, row 34.2 and clauses 5.3.4-5.3.5, 5.6.2]; [2] for colour rendering; [8, Tables 5.4-5.5 and clauses 5.6.1-5.6.2] for energy and controls. *Conservative uniformity bound across a group.}

\subhead{5.1 Illuminance and uniformity}

The background floor receives 617 lx, which is significantly higher than the 100 lx target, because the regular ceiling grid delivers more light than required in circulation areas. To address this, the four circulation zones provide a way to test reduced lighting outputs against the background and nearby task criteria, while group uniformity bounds are calculated using the lowest point divided by the highest surface mean, rounded down.

\subhead{5.2 Glare and screen reflections}

In the baseline, the highest directional UGR is 18.1. To verify luminance limits, the LDT check scales intensity to 4000 lm and divides it by the projected area of the 510 × 510 mm aperture, which gives a maximum sampled luminance of 1334 cd/m² (C44, gamma 65 degrees). This is below the 1500 cd/m² limit set by EN clause 5.9.2 and Table 7 for Case A medium-luminance screens (high state up to 200 cd/m²) tilted up to 15 degrees from vertical [1]. While the check covers sampled angles from 65 to 89 degrees, since projection is undefined at 90 degrees, screen type, tilt, and orientation must be confirm at fit-out.
